% Options for packages loaded elsewhere
\PassOptionsToPackage{unicode}{hyperref}
\PassOptionsToPackage{hyphens}{url}
\PassOptionsToPackage{dvipsnames,svgnames,x11names}{xcolor}
\PassOptionsToPackage{space}{xeCJK}
\documentclass[
  a4paper,
  10pt]{article}
\usepackage{xcolor}
\usepackage{amsmath,amssymb}
\setcounter{secnumdepth}{5}
\usepackage{iftex}
\ifPDFTeX
  \usepackage[T1]{fontenc}
  \usepackage[utf8]{inputenc}
  \usepackage{textcomp} % provide euro and other symbols
\else % if luatex or xetex
  \usepackage{unicode-math} % this also loads fontspec
  \defaultfontfeatures{Scale=MatchLowercase}
  \defaultfontfeatures[\rmfamily]{Ligatures=TeX,Scale=1}
\fi
\usepackage{lmodern}
\ifPDFTeX\else
  % xetex/luatex font selection
  \setmainfont[]{Times New Roman}
  \setmonofont[]{Menlo}
  \setmathfont[]{STIX Two Math}
  \ifXeTeX
    \usepackage{xeCJK}
    \setCJKmainfont[]{Songti SC}
  \fi
  \ifLuaTeX
    \usepackage[]{luatexja-fontspec}
    \setmainjfont[]{Songti SC}
  \fi
\fi
% Use upquote if available, for straight quotes in verbatim environments
\IfFileExists{upquote.sty}{\usepackage{upquote}}{}
\IfFileExists{microtype.sty}{% use microtype if available
  \usepackage[]{microtype}
  \UseMicrotypeSet[protrusion]{basicmath} % disable protrusion for tt fonts
}{}
\setlength{\emergencystretch}{3em} % prevent overfull lines
\providecommand{\tightlist}{%
  \setlength{\itemsep}{0pt}\setlength{\parskip}{0pt}}
\usepackage[a4paper,top=24mm,bottom=25mm,left=23mm,right=23mm,headheight=15pt]{geometry}
\usepackage{amsmath,amssymb,mathtools,bm}
\usepackage{booktabs,longtable,array}
\usepackage{graphicx}
\usepackage{float}
\usepackage{caption}
\usepackage{enumitem}
\usepackage{fancyhdr}
\usepackage{titlesec}
\usepackage{mdframed}
\usepackage{xcolor}
\usepackage{xurl}
\usepackage{hyperref}
\graphicspath{{../../}}

\definecolor{PaperBlue}{HTML}{000000}
\definecolor{PaperRule}{HTML}{B7C6D5}
\definecolor{PaperShade}{HTML}{F4F7FA}
\definecolor{PaperText}{HTML}{000000}

\hypersetup{
  colorlinks=true,
  linkcolor=PaperBlue,
  urlcolor=PaperBlue,
  citecolor=PaperBlue,
  pdfborder={0 0 0}
}

\color{PaperText}
\linespread{1.08}
\setlength{\parindent}{1.5em}
\setlength{\parskip}{0.25em}
\setlength{\emergencystretch}{3em}
\setlist{itemsep=0.2em,topsep=0.45em}
\setcounter{secnumdepth}{3}
\setcounter{tocdepth}{3}

\titleformat{\section}
  {\Large\bfseries\color{PaperBlue}}
  {\thesection}{0.7em}{}
\titleformat{\subsection}
  {\large\bfseries\color{PaperText}}
  {\thesubsection}{0.65em}{}
\titleformat{\subsubsection}
  {\normalsize\bfseries\color{PaperText}}
  {\thesubsubsection}{0.6em}{}
\titlespacing*{\section}{0pt}{2.2ex plus 0.5ex}{0.9ex}
\titlespacing*{\subsection}{0pt}{1.8ex plus 0.4ex}{0.7ex}
\titlespacing*{\subsubsection}{0pt}{1.5ex plus 0.3ex}{0.5ex}

\captionsetup{
  font=small,
  labelfont=bf,
  textfont=it,
  justification=centering,
  skip=6pt
}

\pagestyle{fancy}
\fancyhf{}
\fancyhead[L]{\small\nouppercase{\leftmark}}
\fancyhead[R]{\small Selene's Blog}
\fancyfoot[C]{\small\thepage}
\renewcommand{\headrulewidth}{0.4pt}
\renewcommand{\footrulewidth}{0pt}
\fancypagestyle{plain}{
  \fancyhf{}
  \fancyhead[R]{\small Selene's Blog}
  \fancyfoot[C]{\small\thepage}
  \renewcommand{\headrulewidth}{0.4pt}
}

\renewenvironment{quote}
  {\begin{mdframed}[
      linewidth=0.7pt,
      linecolor=PaperRule,
      backgroundcolor=PaperShade,
      leftline=true,
      rightline=false,
      topline=false,
      bottomline=false,
      innerleftmargin=10pt,
      innerrightmargin=8pt,
      innertopmargin=7pt,
      innerbottommargin=7pt,
      skipabove=8pt,
      skipbelow=10pt
    ]\itshape}
  {\end{mdframed}}

\newenvironment{theorembox}[1]
  {\par\medskip\noindent\textbf{Theorem (#1).}\quad}
  {\par\medskip}

\newenvironment{lemmabox}[1]
  {\par\medskip\noindent\textbf{Lemma #1.}\quad}
  {\par\medskip}

\newenvironment{proofbox}
  {\par\medskip\noindent\textit{Proof.}\quad}
  {\hfill$\square$\par\medskip}

\AtBeginDocument{\allowdisplaybreaks[2]}
\usepackage{bookmark}
\IfFileExists{xurl.sty}{\usepackage{xurl}}{} % add URL line breaks if available
\urlstyle{same}
% fallback for those not using the hyperref driver hyperxmp:
\makeatletter
\@ifundefined{xmpquote}{\newcommand{\xmpquote}[1]{#1}}{}
\makeatother
\hypersetup{
  pdftitle={CS224N \textbar{} GloVe Model},
  pdfauthor={Selene},
  colorlinks=true,
  linkcolor={black},
  filecolor={Maroon},
  citecolor={Blue},
  urlcolor={black},
  pdfcreator={LaTeX via pandoc}}

\title{CS224N \textbar{} GloVe Model}
\author{Selene}
\date{August 6, 2026}

\begin{document}
\maketitle

\section{GloVe Model}\label{glove-model}

\subsection{Comparison with Previous
Models}\label{comparison-with-previous-models}

Previously, we studied the skip-gram model, which mainly learns word
embeddings by making predictions within local context windows. It can
capture complex linguistic patterns beyond word similarity, but it does
not make use of global co-occurrence statistics.

The GloVe model, on the other hand, uses global statistical information
and predicts the probability that word \(j\) appears in the context of
word \(i\) through a least-squares objective.

\subsection{Co-occurrence Matrix}\label{co-occurrence-matrix}

Let \(X\) be the word-word co-occurrence matrix, where \(X_{ij}\)
denotes the number of times word \(j\) appears in the context of word
\(i\)

Let

\[
X_i = \sum_k X_{ik}
\]

denote the total number of times any word \(k\) appears in the context
of word \(i\)

Finally, let

\[
P_{ij}=P(w_j \mid w_i)=\frac{X_{ij}}{X_i}
\]

denote the probability that word \(j\) appears in the context of word
\(i\).

\subsection{Least-Squares Objective
Function}\label{least-squares-objective-function}

In the skip-gram model, we use softmax to compute the probability that
word \(j\) appears in the context of word \(i\):

\[
Q_{ij} = \frac{exp(u_{j}^{T}x_i)}{\sum_{w=1}^{W}exp(u_w^T v_i)}
\]

Training is performed in an online and stochastic manner, using the
following global cross-entropy loss function:

\[
J = -\sum_{j \in corpus}\sum_{j \in context(i)} logQ_{ij}
\]

A significant disadvantage of cross-entropy is that it requires the
distribution \(Q\) to be properly normalized, which requires an
expensive summation over the entire vocabulary. For this reason, we
instead use a least-squares objective and discard the normalization
factors in \(P\) and \(Q\):

\[
\hat{J}=\sum_{i=1}^{W}\sum_{j=1}^{W}X_i\left(\hat{p}_{ij}-\hat{Q}_{ij}\right)^2
\]

where

\[
\hat{P}_{ij}=X_{ij},\qquad
\hat{Q}_{ij}=\exp\left(u_j^{T}v_i\right)
\]

are both unnormalized distributions. Since \(X_{ij}\) is often very
large, optimization becomes difficult. An effective modification is to
minimize the squared error between the logarithms of \(\hat{P}\) and
\(\hat{Q}\):

\[
\hat{J}=\sum_{i=1}^{W}\sum_{j=1}^{W}X_i\left(\log\hat{P}_{ij}-\log\hat{Q}_{ij}\right)^2=\sum_{i=1}^{W}\sum_{j=1}^{W}X_i\left(u_j^{T}v_i-\log X_{ij}\right)^2
\]

\section{Evaluation of Word Vectors}\label{evaluation-of-word-vectors}

\subsection{Intrinsic Evaluation}\label{intrinsic-evaluation}

Intrinsic evaluation of word vectors evaluates a set of word vectors
produced by an embedding method on a specific intermediate subtask. Such
subtasks are usually simple and fast to compute, so they can help us
understand the system that generates the word vectors. Intrinsic
evaluation should usually return a numerical value that indicates how
well these word vectors perform on the evaluation subtask.

\subsection{Extrinsic Evaluation}\label{extrinsic-evaluation}

Extrinsic evaluation of word vectors evaluates a set of word vectors
produced by an embedding method on the real task at hand. Such tasks are
usually more complex and slower to compute. In general, optimizing only
for a poorly performing extrinsic evaluation system cannot determine
exactly which specific subsystem has gone wrong. This also illustrates
the necessity of intrinsic evaluation.

\end{document}
